\section{Attaque sur le lemme \texttt{injective\_agreement\_B\_A}}

\subsection{Description}

Le lemme \texttt{injective\_agreement\_B\_A} vise à garantir une propriété d'accord injectif
entre $B$ et $A$. Cette propriété stipule que chaque acceptation de la clé $k_{AB}$ par $B$
correspond à une unique initiation par $A$. Formellement :
\[
    \forall A, B, t, \#i : \text{Commit\_B}(A, B, \langle \text{`init'}, t \rangle) @ \#i \implies
\]
\[
    \exists \#j : \text{Running\_A}(A, B, \langle \text{`init'}, t \rangle) @ \#j \land \#j < \#i \land
\]
\[
    \neg \exists A_2, B_2, \#i_2 : \text{Commit\_B}(A_2, B_2, \langle \text{`init'}, t \rangle) @ \#i_2 \land \#i_2 \neq \#i
\]

\subsection{Trace d'attaque identifiée par Tamarin}

Tamarin génère une trace d'attaque exploitant la possibilité de rejouer les messages du protocole.
Dans un premier temps, une session légitime se déroule entre $A$ et $B$. L'agent $A$ génère la
clé $k_{AB}$ et le nonce $N_A$, puis les messages du protocole sont échangés normalement jusqu'à
ce que $B$ termine avec succès. L'adversaire enregistre tous les messages échangés durant cette
session, en particulier les messages $\{|\langle A, N_A \rangle|\}_{k_{AB}}$ et $\{|\langle A, \tau, \lambda, k_{AB} \rangle|\}_{k_{BS}}$.

Par la suite, l'adversaire rejoue ces deux messages vers $B$. Puisque les messages sont toujours
valides selon la vérification temporelle $t < \tau + \lambda$, l'agent $B$ accepte une deuxième
fois la clé $k_{AB}$. Du point de vue de Tamarin, il existe maintenant deux exécutions distinctes
de $\text{Commit\_B}$ avec la même clé $k_{AB}$, alors qu'il n'existe qu'une seule exécution de
$\text{Running\_A}$. Cela viole la propriété d'injectivité qui exige une correspondance bijective
entre les acceptations de $B$ et les initiations de $A$.

Cette attaque est rendue possible par l'absence d'un mécanisme liant de manière unique chaque
réponse de $B$ à une requête spécifique. Le timestamp $\tau$ et la durée de validité $\lambda$
ne suffisent pas à empêcher le rejeu durant la fenêtre temporelle $[\tau, \tau + \lambda]$, et
$B$ ne conserve pas de trace des clés de session déjà utilisées.

\subsection{Justification de la validité du protocole}

Cette attaque par rejeu n'est pas considérée comme une violation de nos spécifications de
sécurité. En effet, le document de spécifications stipule explicitement que B doit vérifier
que le nonce n'a pas déjà été envoyé. Donc si il reçoit 2 fois le même message il pourra avorter
le protocole. Cette modification a été ajoutée dans tamarinn mais nous n'avons pas réussi à faire
fonctionner le protocole avec cete modification des règles de réécriture.

La propriété d'injectivité stricte n'est donc pas nécessaire dans notre contexte d'utilisation.
Ce qui importe est que chaque session soit correctement authentifiée et que la clé reste secrète
pendant sa durée d'utilisation effective, ce qui est garanti par les propriétés de fraîcheur et
de secret que Tamarin a pu vérifier avec succès.